% =============================================================
% V. ESTONIAN GROWTH IN COMPARATIVE PERSPECTIVE
% =============================================================
\section{Estonian growth in comparative perspective}
\label{sec:implications}

The series spans three regimes: the capitalist republic of 1919--40, the planned economy of 1940--91 and the market economy restored after 1991. The questions that the literature on catch-up growth asks of each are the same. Gerschenkron's late industrialiser can borrow technology rather than invent it and so grow faster than the leader; the neoclassical model of Solow predicts the same from diminishing returns to capital.\footnote{\autocite[8]{gerschenkron1962}; \autocite[70--1]{solow1956}.} But Abramovitz and David showed that the potential is realised only where a follower possesses the `social capability' to absorb what it borrows, the literacy, institutions and commercial organisation that make imported technology productive, and where the leader's technology is congruent with the follower's endowments.\footnote{\autocite{abramovitz_david1996}.} North, and Acemoglu and Robinson, make the institutional condition explicit: convergence is conditional on institutions that reward productive investment.\footnote{\autocite{north1990}; \autocite{acemoglu_whynationsfail}.} Estonia is an unusually clean test of these propositions, because its social capability was, by the standards of the European periphery, high throughout, while its institutions were changed twice from outside. Tables~\ref{tab:relative}, \ref{tab:interwar}, \ref{tab:trade} and~\ref{tab:structure} provide the comparative evidence.

% tab:relative
\begin{table}[htbp]
\centering\small
\caption[Estonian GDP per capita relative to comparators, selected years]{\textsc{Estonian GDP per capita relative to comparators, selected years}}
\label{tab:relative}
\begin{threeparttable}
\begin{tabular}{lrrrrrrr}
\toprule
& Estonia & \multicolumn{6}{c}{Estonia as a percentage of} \\
\cmidrule(lr){3-8}
Year & (\intlnineteen{}) & Finland & Sweden & Western core & USSR & Lithuania & Latvia \\
\midrule
1913 & 2,058 & 97 & 72 & 51 & 146 & 126 & 87 \\
1919 & 1,600 & 96 & 57 & 48 & 279 & 157 & 108 \\
1929 & 2,113 & 78 & 52 & 45 & 152 & 134 & 90 \\
1938 & 2,685 & 75 & 55 & 52 & 125 & 130 & 87 \\
\midrule
1940 & 2,771 & 86 & 57 & 52 & 129 & 159 & --- \\
1950 & 3,667 & 86 & 54 & 66 & 129 & 160 & --- \\
1968 & 7,348 & 91 & 64 & 70 & 141 & 136 & --- \\
1973 & 8,080 & 73 & 60 & 63 & 133 & 125 & --- \\
1980 & 8,785 & 68 & 59 & 61 & 137 & 127 & --- \\
1989 & 8,958 & 53 & 51 & 52 & 126 & 114 & --- \\
\midrule
1995 & 6,277 & 39 & 35 & 33 & 145 & 121 & 125 \\
2022 & 18,863 & 74 & 64 & 68 & 155 & 97 & 110 \\
\bottomrule
\end{tabular}
\begin{tablenotes}[para,flushleft]
\footnotesize
\textit{Sources:} Estonia: see online appendix. Lithuania: \cite{klimantas2026soviet}; Latvia: \cite{Klimantas2024} (1940--89 not shown). Finland, Sweden, Western core and USSR: \cite{mpd2023} and online appendix. \textit{Notes:} USSR within 1990 borders; 1995 and 2022 refer to the former Soviet Union aggregate.
\end{tablenotes}
\end{threeparttable}
\end{table}


\subsection{Interwar Estonia: recovery, depression and economic nationalism}

The interwar record falls into two phases. The 1920s were a recovery: income per head grew by \gr{1919}{1929} per cent a year but regained its 1913 level only in 1929, and then by barely 3 per cent. The sectors that led, industry, construction and transport, were those that had been most damaged, and the limits of the recovery were set by the permanent loss of the imperial market. Pre-war Estonian industry had been built to serve Russia; the Bolshevik revolution closed that market, and the new republic's attempts to recover a Russian transit trade through Tallinn in the early 1920s came to little.\footnote{\autocite[193--210]{koll1994gap}.} Finland, which had never depended on the Russian market to the same degree and whose forest industries found Western buyers at once, grew by \grx{Finland}{1919}{1929} per cent a year over the same decade (Table~\ref{tab:growth}). In the 1920s, therefore, Estonia's social capability was not in doubt, but the technology it had absorbed before 1914 was no longer congruent with the markets available to it.

The 1930s look different. The Depression was shallow, 8 per cent from peak to trough, and short, and the recovery after 1933 was among the fastest in Europe: Estonian income per head grew by \gr{1929}{1938} per cent a year between 1929 and 1938, when the Western core managed \grx{Western}{1929}{1938} per cent and East-Central Europe \grx{ECE}{1929}{1938} per cent (Table~\ref{tab:growth}). The timing follows the exchange rate. Estonia clung to gold through the banking crisis of 1931--2, suffered the deflation that gold imposed, and recovered only after it devalued and pegged to sterling in June 1933, a pattern consistent with that which Eichengreen and Sachs identified for Europe as a whole and that the Nordic economies displayed.\footnote{\autocite[74]{bernankejames01}; \autocite{RefWorks:doc:5ea847d2e4b0aa025739fe41}; \autocite{grytten_nordicdepression}; \autocite[287--99]{kuokstis2016}.} The content of the recovery was the policy that Kõll and Valge have called economic nationalism: state investment in oil shale, textiles and food processing, protection of the home market and the organisation of agricultural exports through co-operatives in imitation of Denmark.\footnote{\autocite[11--29]{koll1998}; \autocite{norkus_denmarkisation}; \autocite[118--35]{klesment2005}.} The sectoral indices confirm that the 1930s acceleration was industrial and urban: construction grew by \sgr{Construction}{1929}{1938} per cent a year, industry by \sgr{Industrytotal}{1929}{1938} per cent and the housing stock by \sgr{Housing}{1929}{1938} per cent, while agriculture grew by \sgr{Primarysectortotal}{1929}{1938} per cent.

Table~\ref{tab:interwar} places the result in its regional setting. In 1938 Estonia's income per head, \ee{1938}~\intlnineteen{}, was above that of Lithuania, Poland and Hungary, slightly below Czechoslovakia's and well below Latvia's, Finland's and the Nordic average. Its growth over 1929--38 was exceeded in the region only by its two Baltic neighbours, by Finland and by the Soviet Union, whose own figures for the period are the subject of the debate reviewed in Section~\ref{sec:method}, and was far above that of the Western core and of East-Central Europe. Table~\ref{tab:structure} shows why the result should be read as a success and why it was not a greater one. Estonia's adult literacy, around 92 per cent, was second only to Latvia's in the European periphery and above Finland's; its university enrolment per head was among the highest in Europe; its land was in the hands of smallholders after the reform of 1919. These are the ingredients of social capability. Yet 63 per cent of the labour force still worked in agriculture in 1934, a share as high as Finland's, though below Latvia's, and agriculture produced some 38 per cent of GDP in 1937 (Table~\ref{tab:structure}). An economy of that structure, open to trade (Table~\ref{tab:trade}) but exporting butter, bacon, timber and flax, could grow at the pace of its agricultural productivity and of the markets Britain and Germany allowed it, and no faster; the industrial drive of 1934--9 was a response to that constraint, not an escape from it.

% tab:interwar
\begin{table}[htbp]
\centering\small
\caption[GDP per capita in Estonia and European comparators, 1913--38]{\textsc{GDP per capita in Estonia and European comparators, 1913--38}}
\label{tab:interwar}
\begin{threeparttable}
\begin{tabular}{lrrrrrrrr}
\toprule
& \multicolumn{5}{c}{\intlnineteen{}} & \multicolumn{3}{c}{Growth, \% p.a.} \\
\cmidrule(lr){2-6}\cmidrule(lr){7-9}
Country & 1913 & 1919 & 1929 & 1932 & 1938 & 1919--29 & 1929--38 & 1913--38 \\
\midrule
Denmark & 3,912 & 3,860 & 5,075 & 5,169 & 5,762 & 2.8 & 1.4 & 1.6 \\
Sweden & 2,874 & 2,825 & 4,063 & 3,968 & 4,901 & 3.7 & 2.1 & 2.2 \\
\textbf{Finland} & 2,111 & 1,658 & 2,717 & 2,550 & 3,589 & 5.1 & 3.1 & 2.1 \\
Austria & 3,465 & 2,259 & 3,699 & 2,940 & 3,559 & 5.1 & $-$0.4 & 0.1 \\
Latvia & 2,361 & 1,481 & 2,351 & 2,286 & 3,099 & 4.7 & 3.1 & 1.1 \\
Ireland & 2,736 & 2,595 & 2,824 & 2,885 & 3,052 & 0.8 & 0.9 & 0.4 \\
Czechoslovakia & 2,096 & 1,588 & 3,042 & 2,680 & 2,803 & 6.7 & $-$0.9 & 1.2 \\
\textbf{Estonia} & 2,058 & 1,600 & 2,113 & 2,105 & 2,685 & 2.8 & 2.7 & 1.1 \\
Greece & 1,455 & 1,274 & 2,342 & 2,289 & 2,677 & 6.3 & 1.5 & 2.5 \\
Hungary & 2,098 & 1,484 & 2,476 & 2,192 & 2,655 & 5.3 & 0.8 & 0.9 \\
Poland & 1,739 & 1,348 & 2,117 & 1,658 & 2,182 & 4.6 & 0.3 & 0.9 \\
USSR & 1,414 & 573 & 1,386 & 1,439 & 2,150 & 9.2 & 5.0 & 1.7 \\
Lithuania & 1,632 & 1,017 & 1,577 & 1,702 & 2,057 & 4.5 & 3.0 & 0.9 \\
\bottomrule
\end{tabular}
\begin{tablenotes}[para,flushleft]
\footnotesize
\textit{Sources:} Estonia: see online appendix. Latvia: 1913, \cite[603]{norkusmark2021}; 1919--38, \cite[316]{klimantas2025thesis}, table E.1, the series published for 1920--39 in \cite{Klimantas2024}, whose values for 1929 and 1932 (2,349 and 2,288) differ from those used here by less than 0.1 per cent. Lithuania: \cite{klimantas2023lithuania}. Other countries: \cite{mpd2023}. \textit{Notes:} countries ordered by 1938 level.
\end{tablenotes}
\end{threeparttable}
\end{table}


% tab:trade
\begin{table}[htbp]
\centering\small
\caption[Trade openness (exports plus imports as a percentage of GDP), period averages]{\textsc{Trade openness (exports plus imports as a percentage of GDP), period averages}}
\label{tab:trade}
\begin{threeparttable}
\begin{tabular}{lrrr}
\toprule
Country & 1924--29 & 1930--34 & 1935--38 \\
\midrule
\textbf{Estonia} & 33.4 & 25.6 & 31.5 \\
Latvia & 38.1 & 26.9 & 21.8 \\
Lithuania & 19.9 & 21.8 & 21.6 \\
\textbf{Finland} & 41.2 & 26.8 & 34.3 \\
Denmark & 51.8 & 35.2 & 36.3 \\
Norway & 56.1 & 35.6 & 41.1 \\
Sweden & 36.8 & 23.9 & 29.3 \\
Poland & 13.3 & 10.3 & 11.8 \\
Czechoslovakia & 36.8 & 23.6 & 26.5 \\
Hungary & 21.7 & 14.1 & 16.8 \\
Germany & 25.0 & 16.9 & 15.6 \\
United Kingdom & 36.9 & 23.3 & 25.6 \\
\bottomrule
\end{tabular}
\begin{tablenotes}[para,flushleft]
\footnotesize
\textit{Sources:} see the online appendix for the compilation of trade and GDP series. Baltic exports and imports: \cite{lidaexportsee}; \cite{lidaimportsee}; \cite{lidaexportslv}; \cite{lidaimportslv}; \cite{lidaexp}; \cite{lidaimp}; \cite[123, 161]{balticstates1938}. Other countries: \cite{worldtrade1949}. \textit{Notes:} unweighted averages of annual ratios; missing years omitted.
\end{tablenotes}
\end{threeparttable}
\end{table}


% tab:structure
\begin{table}[htbp]
\centering\footnotesize\setlength{\tabcolsep}{4pt}
\caption[Social capability and economic structure around 1930: Estonia, Finland and European comparators]{\textsc{Social capability and economic structure around 1930: Estonia, Finland and European comparators}}
\label{tab:structure}
\begin{threeparttable}
\begin{tabular}{lrrrr}
\toprule
Country & Literacy & Agriculture & University & Agriculture \\
& (\% of adults) & (\% of employment) & students (per 1,000) & (\% of GDP) \\
& c.\,1930 & c.\,1930 & 1930 & c.\,1937 \\
\midrule
Sweden & 99.6 & 36.0 & 1.6 & 13.0 \\
United Kingdom & 97.5 & 6.0 & 3.1 & 4.0 \\
Germany & 97.5 & 29.0 & --- & 16.0 \\
Denmark & 97.5 & 35.3 & 2.7 & 18.0 \\
Norway & 97.5 & 35.5 & 2.2 & 13.0 \\
Austria & 97.5 & 39.0 & 3.2 & 14.0 \\
Ireland & 97.5 & 48.5 & 1.5 & 24.8 \\
Czechoslovakia & 96.0 & 37.5 & --- & 24.0 \\
Latvia & 93.4 & 67.2 & 4.9 & 39.6 \\
\textbf{Estonia} & 92.0 & 63.0 & 3.5 & 37.6 \\
Hungary & 90.0 & 54.2 & --- & 31.0 \\
\textbf{Finland} & 84.1 & 64.6 & 2.0 & 32.0 \\
Italy & 78.4 & 48.2 & 1.1 & 29.0 \\
Lithuania & 77.0 & 75.2 & 1.8 & 43.8 \\
Poland & 73.0 & 65.9 & --- & --- \\
Bulgaria & 69.0 & 71.8 & --- & 52.1 \\
Spain & 68.0 & 51.9 & 1.7 & 22.0 \\
Greece & 58.1 & 56.6 & 1.1 & 54.4 \\
Romania & 57.1 & 77.2 & --- & --- \\
Yugoslavia & 55.0 & 79.7 & --- & 41.0 \\
USSR & 50.0 & 73.0 & --- & 29.0 \\
Portugal & 39.8 & 48.0 & 0.8 & 30.0 \\
\bottomrule
\end{tabular}
\begin{tablenotes}[para,flushleft]
\footnotesize
\textit{Sources:} literacy: \cite{cronin_literacy}; \cite[246--77]{kirk_literacy}; \cite[528]{bostonencikl_rast}. Agricultural employment: \cite[181]{morysinterwarinmorys2021}; \cite{mitchelstat}; \cite[395, 409, 418]{sov_labour1939}. University students and agricultural shares of GDP: see online appendix, section A.7. \textit{Notes:} countries ordered by literacy; a literacy value of 97.5 denotes near-universal literacy where the source gives no census figure. `---': not available.
\end{tablenotes}
\end{threeparttable}
\end{table}


\subsection{Under Soviet rule: collectivisation, industrialisation and stagnation}

The Soviet period opens with a decade of destruction. Income per head fell by about a third between 1940 and 1944 on the present estimates, and the population by a fifth between 1939 and 1945 through war deaths, deportation and flight.\footnote{\autocite[265]{taagepera_coll01}; \autocite[12]{deportations_ee}.} The recovery of industry and construction after 1945 was rapid, as Soviet reconstruction everywhere was: both regained their 1938 level by 1948, and industry grew by around 20 and construction by around 14 per cent a year between 1945 and 1953 (Tables~\ref{tab:sectoral} and~\ref{tab:appsec1}). Agriculture tells a different story, and one that the sectoral indices allow to be told with some precision. Livestock output, halved by the war, was still 38 per cent below its 1938 level in 1950, the year after collectivisation; crop output recovered sooner, but the agricultural sector as a whole stood at 60--65 per cent of its 1938 level in 1945--6, fluctuated between 77 and 93 per cent of it in 1947--57, and exceeded it only in \sregain{Primarysectortotal}. The index shows a shallower trough than Mertelsmann's archival estimate, which puts official harvests 10--15 per cent above reality and output after collectivisation at perhaps 40 per cent of the 1939 level, but the two agree that the sector did not recover for a decade; Mertelsmann also finds that the collective and state farms produced little more than a third of agricultural output in the early 1950s, the rest coming from the private plots of their members, which still supplied 28 per cent of output in 1960.\footnote{\autocite[156--8]{olaf_stalin}; \autocite[79--109]{feest2014}; \autocite[257]{klesment2009}; \autocite[160]{hansen1994}.} Nothing comparable to the Soviet famine of 1932--3 occurred, probably because the Baltic collectivisation of 1949 did not coincide with a procurement crisis and because, as Harrison has argued for Soviet agriculture generally, the system had by then learned to extract the surplus without destroying the herds.\footnote{\autocite[192--208]{harrison1996}.} The cost was a lost decade and a permanent shift of effort towards animal husbandry, in which Estonia came to lead the Union in output per head, while crop output peaked in \speak{Crops} and never grew again.\footnote{\autocite[234]{years_of_dependence}.}

Between 1950 and 1968 Estonia experienced what Gerschenkron would have recognised as a classic burst of state-led industrialisation. Total GDP grew by \tgr{1950}{1968} per cent a year and income per head by \gr{1950}{1968} per cent, faster than in Finland, the Nordic countries or the Western core (Table~\ref{tab:growth}). The engine was investment directed from Moscow into three complexes. The first was the oil-shale basin of the north-east, which fed two of the largest thermal power stations in the Union, completed in 1965 and 1973; electricity output rose from 2 billion kilowatt-hours in 1960 to 12 billion in 1970, of which Estonia itself consumed a third. The second was cotton textiles, which Moscow expanded in Estonia rather than in cotton-growing Uzbekistan because Estonia had the skilled labour and the mills. The third was the machine-building and electronics plants that supplied the Union market.\footnote{\autocite[239]{years_of_dependence}; \autocite[41--7]{ahderajasalu1992}; \autocite[248--58]{pihlamagi2013}.} The sovnarkhoz reform of 1957--65, which briefly placed industrial management in Tallinn, saw industrial output double in seven years, with 70 per cent of investment going to industry and the share of light industry in the total deliberately reduced.\footnote{\autocite[248--52]{pihlamagi2013}; \autocite[67--70]{purre1965}.} Heavy industry grew by \sgr{Heavyindustries}{1950}{1968} per cent a year, textiles by \sgr{Textiles}{1950}{1968} per cent and construction by \sgr{Construction}{1950}{1968} per cent. The labour came partly from a countryside that collectivisation had released and partly from immigration: Estonia's population grew by half between 1950 and 1989 while the Estonian share of it fell from more than 90 to little more than 60 per cent.\footnote{\autocite[158--217]{case_estonia}.}

How did this growth compare with that of the Union into which Estonia had been absorbed? Table~\ref{tab:relative} gives the answer, and it is not the one that either Soviet or nationalist accounts have assumed. Estonia entered the Union in 1940 with income per head \pctussr{1940} per cent of the Soviet average; it stood at \pctussr{1968} per cent in 1968 and at \pctussr{1989} per cent in 1989. The gap neither closed nor widened appreciably over half a century: Estonia grew with the Union, somewhat faster in the 1950s and 1960s and somewhat slower thereafter, and remained one of its two or three richest republics throughout. The result matches the Soviet-era estimates of relative republican income, which placed Estonia at 137--145 per cent of the Union average in 1960--70 by net material product and at 134 per cent in 1970 on Koropeckyj's reworking, and it matches the CIA's finding that Estonian national income per head grew by 6.3 per cent a year in 1961--9 against 5.5 per cent for the Union.\footnote{\autocite[169, 174]{schroeder1973}; \autocite[64--71]{bs_livingstandards1976}; \autocite[8--10]{cia1972}; \autocite[35--8]{bs_livingstandards1988}.} It also bears on the question whether Moscow bought the loyalty of its western borderlands with investment. Investment per head in Estonia was above the Union average in the 1960s, second only to Kazakhstan, but the republic's share of investment in its own national income was below the Union average. Van Selm and Dölle's reconstruction of inter-republican transfers finds Estonia's net balance close to zero from 1966 to 1978, becoming a net recipient of 4--6 per cent of national income only after 1979 and of 13 per cent in 1988, when the Union economy was already unravelling.\footnote{\autocite[22--9]{cia1972}; \autocite[53]{adirim1988}; \autocite[136--40]{vanselmdolle1993}.} For the years after 1966 for which transfers can be measured, Estonia's high income under Soviet rule was earned rather than transferred; and over the period as a whole it was inherited rather than created, since the republic was richer than the Union when it was annexed and remained so in roughly the same proportion.

The deceleration after 1968 was general and, by the late 1970s, complete. Income per head grew by \gr{1968}{1989} per cent a year over 1968--89 and by \gr{1973}{1989} per cent after 1973, and the sectoral peaks cluster in a single decade: food processing in \speak{Food}, construction in \speak{Construction}, electricity and oil-shale output in 1979--80, heavy industry in \speak{Heavyindustries}.\footnote{Electricity peaked in 1979 and fell by 9 per cent to 1985; oil-shale mining levelled off after the arrival of natural gas (\autocite[239]{years_of_dependence}; \autocite{taagepera2018}).} The sources point to three causes. The Union-wide slowdown in investment growth that the tenth and eleventh five-year plans deliberately imposed fell on a republic whose growth had been investment-led, and whose own labour reserves were exhausted: contemporaries noted that further industrial expansion in Estonia required either immigration or capital deepening, and the planners increasingly withheld both.\footnote{\autocite[128--63]{hanson2003}; \autocite[164--76]{kolde1978}; \autocite[441]{dienes1972}.} The oil-shale complex, built on a resource that was expensive to mine and polluting to burn, lost its rationale once Siberian gas reached the Baltic.\footnote{\autocite[188--9]{olaf_stalin}.} And the gains in reported output after 1975 were increasingly nominal: Finnish and Estonian economists who compared the two economies in 1992 concluded that the rapid growth of Estonian industry to 1975 reflected real investment, but that `the further growth of production can be explained mainly by the change in prices'.\footnote{\autocite[39--40]{ahderajasalu1992}.} The physical indices used here confirm that judgement: industrial output grew by \sgr{Industrytotal}{1968}{1989} per cent a year after 1968, and the branches that continued to grow, textiles and other light industry, did so slowly, while cotton cloth output was flat from 1970.\footnote{\autocite[41]{ahderajasalu1992}.} Two further irregularities in the indices deserve mention. Construction, which had grown by about 9 per cent a year since 1950, peaked in \speak{Construction} and stagnated for the rest of the Soviet period as the Union-wide investment squeeze took hold, and crop output fell by a quarter in each of the poor harvests of 1978 and 1981 that affected the Union as a whole; and the housing stock, which had grown by 2 per cent a year since the 1950s, stopped growing in 1986 and declined thereafter, as completions fell from 13,500 dwellings a year in 1985 to 5,200 in 1991.\footnote{\autocite[134--50]{hanson2003}; \autocite[60, 66]{ahderajasalu1992}; \autocite[216]{wb1993estonia}.}

Read through the lens of growth theory, the Soviet episode is coherent. The catch-up of 1950--68 was what the models predict when a capital-poor economy with high social capability is made to invest: Estonia's literate and urbanised labour force absorbed Soviet technology quickly, and the mass-production methods of mid-century were congruent with its endowments. The stagnation that followed is equally what the models predict when the institutional condition fails. Extensive growth exhausted its inputs, and the system could not replace them with the productivity growth that flexible production, foreign technology and competition were delivering across the Gulf: the share of services in Estonian GDP was 42 per cent in 1990, barely above its interwar level, against about 60 per cent in Finland.\footnote{Table~\ref{tab:benchmarks}; \autocite[11--12]{ahderajasalu1992}; \autocite{klimantas2026soviet}.} Estonia's income relative to the Western core, which had risen from \pctwc{1950} per cent in 1950 to \pctwc{1968} per cent in 1968, fell back to \pctwc{1989} per cent by 1989, almost exactly its level of 1913. Two decades of Soviet growth had been cancelled by two decades of Soviet stagnation.

\subsection{Transition and recovery}

The transition compressed into four years a fall in income per head of \eeTROUGHpct{} per cent, deeper than the Depression and comparable only to the two world wars. The sectoral indices to 1991 show the collapse beginning in Soviet-era industry and construction, which contracted at \sgr{Industrytotal}{1989}{1991} and \sgr{Construction}{1989}{1991} per cent a year even before the rouble was abandoned. After 1991 the reconstruction follows the modern series, which records the loss of eastern markets and cheap Soviet inputs, the currency board of 1992 and the stabilisation that followed.\footnote{\autocite[1--8]{wb1993estonia}; \autocite[164--6]{norkus_restorations}; \autocite[482--509]{shteinbuka1993}.} By comparison with the other Soviet republics the fall was contained, probably because Estonia had less heavy industry to lose and reoriented trade towards Finland and Sweden faster than any other successor state.\footnote{\autocite[22]{transition_greenes}; \autocite[443--4]{transition_progress}.} Recovery from 1995 at \gr{1995}{2022} per cent a year carried Estonia from \pctwc{1995} per cent of the Western core in 1995 to \pctwc{2022} per cent in 2022, close to the \pctwc{1968} per cent reached in 1968 at the height of Soviet growth. Whether the catch-up of the market era will continue or stall in the middle-income trap that Staehr has described is a question for another paper; what the long series shows is that Estonia converged on Western Europe under capitalism twice, in the 1930s and after 1995, and under planning once, briefly, before diverging again.\footnote{\autocite[274--80]{staehr2015}.}
